\documentclass[main.tex]{subfiles}

\begin{document}

\section{Theory Appendix: Proofs and Details}\label{sec:appendix-theory}

The sufficient statistic rests on an envelope theorem for the consumer: when interchange fees change, consumers care only about the change in retail prices on the goods they buy and the change in the reward on the method they use.
The hard part of the proof is to show that those equilibrium price changes are well approximated by functions of equilibrium quantities we can measure at the merchant---the composition of payments and the fee schedule.

\subsection{Setup and Notation}\label{subsec:setup-notation}

We prove Theorem~\ref{thm:sufficient-stat} in the sufficient-statistic environment of Section~\ref{sec:redistribution}.
Consumers, merchants, and networks behave as in that section: type-$k$ consumers maximize $U_k$ subject to $\sum_j p_j q_{jk}\leq 1+f_k$; interchange fees follow~\eqref{eq:interchange-fee-model}; retail prices follow~\eqref{eq:retail-price-pass}; and rewards follow~\eqref{eq:reward-pass-through}.
We use stars for equilibrium values and index payment methods by $k$ and $l$.

For the proofs it is convenient to write each merchant's payment-weighted average fee as $\hat\tau_j\equiv\sum_m\mu_{jm}\tau_{jm}$, so~\eqref{eq:retail-price-pass} reads $\log p_j=\log\bar p_j+\rho^P\hat\tau_j$.
Revenue weights $\omega_j$ and the operator $\mathbb{E}_R[\cdot]\equiv\sum_j\omega_j(\cdot)$ are as in Theorem~\ref{thm:sufficient-stat}; let $\tau_{\max}\equiv\max_{j,k}|\tau_{jk}|$.
Rewards inflate each cohort's effective budget to $1+f_k$, so aggregate expenditure is $1+O(\tau_{\max})$ (Assumption~\ref{assump:positive-costs}); in all first-order expressions below we treat aggregate expenditure as one and absorb the $O(\tau_{\max})$ correction into the stated remainders.
The proof uses the following additional notation.

\begin{center}
\begin{tabular}{lll}
\toprule
\textbf{Symbol} & \textbf{Dimension} & \textbf{Meaning} \\
\midrule
$\hat\tau_j$ & scalar & merchant $j$'s payment-weighted average fee, $\sum_m\mu_{jm}\tau_{jm}$ \\
$\tau_{jk}-\hat\tau_j$ & scalar & price-side ``excess fee'': method $k$'s deviation from merchant $j$'s average \\
$s^*_{jk}$ & scalar & cohort $k$'s spending share at merchant $j$, $p_jq^*_{jk}/\sum_{j'}p_{j'}q^*_{j'k}$ \\
$\bar\tau^k$ & scalar & cohort $k$'s spending-weighted average fee, $\sum_j s^*_{jk}\tau_{jk}$ \\
$E^{k}$ & $\mathbb{R}^{J\times J}$ & price elasticities of cohort $k$: $E^k_{ji}=\partial\log q^k_j/\partial\log p_i$ \\
$G^{k}$ & $\mathbb{R}^{J\times K}$ & income semi-elasticities of cohort $k$: $G^k_{jm}=\partial\log q^k_j/\partial f_m$ \\
\bottomrule
\end{tabular}
\end{center}

\subsection{Assumptions}

Let $E^k_{ji}\equiv\partial\log q^k_j/\partial\log p_i$ denote the cross-price elasticity of demand among consumers of type $k$, and let $E$ stack these matrices across types.

\begin{assumption}\label{assump:dominant-diagonal}
The elasticities $E^k_{ji}$ are finite and satisfy a dominant diagonal condition: there exist $M,m>0$ such that for each $k$ and $j$,
\[
|E^k_{jj}|<M\quad\text{and}\quad\sum_{i\neq j}|E^k_{ji}|<m\,|E^k_{jj}|.
\]
\end{assumption}

Own-price elasticities are bounded above and dominate the sum of cross-price elasticities for the same good.
CES demand with $\sigma>1$ satisfies the condition with $M=\sigma$ and $m=1$.

\begin{assumption}[Rewards of Fee Order]\label{assump:positive-costs}
Observed rewards are the same order of magnitude as interchange fees, $f_k=O(\tau_{\max})$.
\end{assumption}

The reward pass-through rule itself is part of the environment in Section~\ref{sec:redistribution}.
The only additional restriction used here is the size condition in Assumption~\ref{assump:positive-costs}: rewards run about $1$--$2\%$ of spend in the data, the same order as interchange fees.

\begin{assumption}[Bounded Income Semi-Elasticities]\label{assump:income-elasticity}
The demand of a consumer of type $k$ admits finite income semi-elasticities $G^k_{jm}\equiv\partial\log q^k_j/\partial f_m$, with row sums bounded uniformly: $\sum_m|G^k_{jm}|\leq L$ for some finite $L$ and for all $j$ and $k$.
\end{assumption}

Homothetic demand is a sufficient special case, with $G^k_{jm}=\mathbbm{1}\{m=k\}/(1+f_k)$ and hence $L=1$.
The three demand bounds $M$, $m$, and $L$ are the only primitives that enter the theorem: $M$ controls how demand responds to retail prices, $m$ controls how demand responds to cross-price elasticities, and $L$ controls how demand responds to rewards (an income change for the cohort).
All three reappear in the contraction threshold of Step~C.

\subsection{Proof of Theorem~\ref{thm:sufficient-stat}}\label{subsec:proof-theorem1}

The proof has four steps.
Step~A applies the envelope theorem to reduce welfare to a price term and a reward term.
Step~B evaluates the reward term from~\eqref{eq:reward-pass-through}.
Step~C gives an explicit bound showing that equilibrium feedback is small relative to the direct price and reward effects.
Step~D assembles the pieces.

\paragraph{Step A: Envelope theorem.}
Roy's identity applied to the consumption problem of a consumer in cohort $k$ gives, using the spending shares $s^*_{jk}$ and total expenditure $1+O(\tau_{\max})$,
\begin{equation}
\frac{1}{\lambda_k}\frac{dV_k}{dx_l}=-\sum_j s^*_{jk}\,\frac{d\log p_j}{dx_l}+\frac{df_k}{dx_l}+O(\tau_{\max}).
\label{eq:cohort-roy}
\end{equation}
Because the consumer was already optimizing, reallocating her basket has no first-order welfare effect; she is hit only through retail prices on the goods she buys, weighted by her spending, and through the reward on the method she uses.

We carry out the proof in these cohort spending shares---the natural weights for the envelope and reward steps---and translate the result into the observable revenue-weighted moments of Theorem~\ref{thm:sufficient-stat} only at the end (Step~D).
The bridge is a spending identity: cohort $k$ spends $\mu_k\,p_j q^*_{jk}$ at merchant $j$, while merchant $j$ draws the share $\omega_j\,\mu_{jk}$ of economy-wide revenue from method $k$; equating the two accounting conventions gives
\begin{equation}
\mu_k\,s^*_{jk}=\omega_j\,\mu_{jk}+O(\tau_{\max}),
\label{eq:spending-identity}
\end{equation}
where the remainder collects reward-induced deviations from unit aggregate expenditure.
Equivalently, a cohort spending-weighted sum is a revenue-weighted moment, $\sum_j s^*_{jk}(\cdot)=\mu_k^{-1}\,\mathbb{E}_R[\mu_{jk}(\cdot)]+O(\tau_{\max})$.
It remains to characterize the reward response $df_k/dx_l$ (Step~B) and the price response $d\log p_j/dx_l$ (Step~C).

\paragraph{Step B: Reward response decomposition.}
The reward rule~\eqref{eq:reward-pass-through} is itself a level identity in spending shares, $f_k=\bar f_k+\rho^R\,\bar\tau^k$ with $\bar\tau^k=\sum_j s^*_{jk}\tau_{jk}$: rewards on method $k$ equal the fixed baseline $\bar f_k$ plus the marginal pass-through $\rho^R$ times the average fee that cohort $k$'s spending encounters.

Differentiate in $x_l$.
By~\eqref{eq:reward-pass-through}, the baseline $\bar f_k$ collects components that do not move with the fee schedule, so $d\bar f_k/dx_l=0$; also $d\tau_{jk}/dx_l=\mathbbm{1}\{k=l\}\,\partial\tau_{jl}/\partial x_l$, so
\begin{equation}
\frac{df_k}{dx_l}=\rho^R\!\left(\mathbbm{1}\{k=l\}\sum_j s^*_{jk}\,\frac{\partial\tau_{jl}}{\partial x_l}+\sum_j\tau_{jk}\,\frac{ds^*_{jk}}{dx_l}\right).
\label{eq:reward-decomp}
\end{equation}
The first term is the direct network-balance rebate: every dollar the network collects from the rate hike on method $l$ is rebated to method-$l$ users.
The second term is the reward feedback from the movement of cohort $k$'s spending shares.
Step~C bounds this feedback term.
Because the baseline $\bar f_k$ is fee-invariant, it drops out of the derivative and leaves the reward response---and hence Theorem~\ref{thm:sufficient-stat}---unchanged; the level of network costs and issuer margins never enters.

\paragraph{Step C: Equilibrium feedback (explicit bound).}
Differentiating~\eqref{eq:retail-price-pass}---equivalently, $\log p_j=\log\bar p_j+\rho^P\hat\tau_j$---and recentering the composition shift with $\sum_m d\mu_{jm}/dx_l=0$ (shares sum to one) gives
\begin{equation}
\frac{d\log p_j}{dx_l}=\rho^P\!\left[\underbrace{\mu_{jl}\,\frac{\partial\tau_{jl}}{\partial x_l}}_{\text{direct effect}}+\underbrace{\sum_m(\tau_{jm}-\hat\tau_j)\,\frac{d\mu_{jm}}{dx_l}}_{\text{composition feedback}}\right].
\label{eq:price-decomp}
\end{equation}
The direct term is the mechanical pass-through of the fee into prices.
The composition feedback is small because it multiplies share movements by centered fee gaps, each of which is at most order $\tau_{\max}$.

Let
\[
z=\begin{pmatrix}d\log p\\[2pt] df\end{pmatrix},\qquad
z_0=\begin{pmatrix}z^p_0\\[2pt] z^f_0\end{pmatrix},
\]
where
\[
z^{p}_{0,j}=\rho^P\,\mu_{jl}\,\frac{\partial\tau_{jl}}{\partial x_l},
\qquad
z^{f}_{0,k}=\rho^R\,\mathbbm{1}\{k=l\}\sum_j s^*_{jl}\,\frac{\partial\tau_{jl}}{\partial x_l}.
\]
Thus $z_0$ is the response holding payment shares and spending shares fixed.
Let
\begin{equation}
B_l\equiv\max_j\left|\frac{\partial\tau_{jl}}{\partial x_l}\right|,
\qquad
H\equiv(1+m)M+L,
\qquad
C\equiv2(1+H).
\label{eq:direct-effect}
\end{equation}
Because $\rho^P,\rho^R\leq1$, $\mu_{jl}\leq1$, and $\sum_j s^*_{jl}=1$, $\|z_0\|_\infty\leq B_l$.

It remains to bound $z-z_0$.
Demand responds to prices and rewards according to
\[
\frac{d\log q^k_j}{dx_l}
=\sum_i E^k_{ji}\frac{d\log p_i}{dx_l}
+\sum_n G^k_{jn}\frac{df_n}{dx_l}.
\]
The row-sum bounds in Assumptions~\ref{assump:dominant-diagonal} and~\ref{assump:income-elasticity} imply, for every $j$ and $k$,
\begin{equation}
\left|\frac{d\log q^k_j}{dx_l}\right|
\leq H\|z\|_\infty .
\label{eq:demand-response-bound}
\end{equation}

For prices, the within-merchant payment share $\mu_{jm}=\mu_m q^*_{jm}/\sum_n\mu_n q^*_{jn}$ nets out the common price $p_j$, so the quotient rule gives
\[
\frac{d\log\mu_{jm}}{dx_l}
=\frac{d\log q^m_j}{dx_l}
-\sum_n\mu_{jn}\frac{d\log q^n_j}{dx_l}.
\]
Since $\sum_m(\tau_{jm}-\hat\tau_j)\mu_{jm}=0$, the average term drops out of the recentered composition feedback:
\[
\sum_m(\tau_{jm}-\hat\tau_j)\frac{d\mu_{jm}}{dx_l}
=\sum_m(\tau_{jm}-\hat\tau_j)\mu_{jm}\frac{d\log q^m_j}{dx_l}.
\]
The centered merchant fee gap satisfies
\[
\sum_m|(\tau_{jm}-\hat\tau_j)\mu_{jm}|\leq 2\tau_{\max}.
\]
Combining this inequality with~\eqref{eq:demand-response-bound},
\begin{equation}
\left|
\frac{d\log p_j}{dx_l}
-\rho^P\,\mu_{jl}\frac{\partial\tau_{jl}}{\partial x_l}
\right|
\leq 2H\tau_{\max}\|z\|_\infty .
\label{eq:price-feedback-bound}
\end{equation}

For rewards, the feedback comes from the spending-share term left in Step~B.
The spending share $s^*_{jk}=p_jq^*_{jk}/\sum_{j'}p_{j'}q^*_{j'k}$ has derivative
\[
\frac{d\log s^*_{jk}}{dx_l}
=\left(\frac{d\log p_j}{dx_l}+\frac{d\log q^k_j}{dx_l}\right)
-\sum_{j'}s^*_{j'k}\left(\frac{d\log p_{j'}}{dx_l}+\frac{d\log q^k_{j'}}{dx_l}\right).
\]
Recentering by $\bar\tau^k$ makes the average term vanish:
\[
\sum_j\tau_{jk}\frac{ds^*_{jk}}{dx_l}
=\sum_j s^*_{jk}(\tau_{jk}-\bar\tau^k)
\left(\frac{d\log p_j}{dx_l}+\frac{d\log q^k_j}{dx_l}\right).
\]
The centered cohort fee gap satisfies
\[
\sum_j|s^*_{jk}(\tau_{jk}-\bar\tau^k)|\leq2\tau_{\max}.
\]
Therefore
\begin{equation}
\left|
\frac{df_k}{dx_l}
-\rho^R\,\mathbbm{1}\{k=l\}\sum_j s^*_{jl}\frac{\partial\tau_{jl}}{\partial x_l}
\right|
\leq 2(1+H)\tau_{\max}\|z\|_\infty .
\label{eq:reward-feedback-bound}
\end{equation}

Equations~\eqref{eq:price-feedback-bound} and~\eqref{eq:reward-feedback-bound} together imply
\begin{equation}
\|z-z_0\|_\infty\leq C\tau_{\max}\|z\|_\infty,
\qquad
C\equiv2\bigl[1+(1+m)M+L\bigr].
\label{eq:contraction-bound}
\end{equation}
Set $\bar\tau\equiv1/C$.
The same calculation applies to the difference between any two candidate responses, so whenever $\tau_{\max}<\bar\tau$, the feedback map is a contraction.
For the equilibrium response,
\[
\|z\|_\infty
\leq \|z_0\|_\infty+\|z-z_0\|_\infty
\leq B_l+C\tau_{\max}\|z\|_\infty,
\]
so
\[
\|z\|_\infty\leq\frac{B_l}{1-C\tau_{\max}}
\qquad\text{and}\qquad
\|z-z_0\|_\infty
\leq\frac{C\tau_{\max}}{1-C\tau_{\max}}B_l .
\]
The coordinate bounds are therefore
\begin{align}
\left|
\frac{d\log p_j}{dx_l}
-\rho^P\,\mu_{jl}\frac{\partial\tau_{jl}}{\partial x_l}
\right|
&\leq
\frac{2H\tau_{\max}}{1-C\tau_{\max}}B_l,
\label{eq:price-step-explicit}\\
\left|
\frac{df_k}{dx_l}
-\rho^R\,\mathbbm{1}\{k=l\}\sum_j s^*_{jl}\frac{\partial\tau_{jl}}{\partial x_l}
\right|
&\leq
\frac{2(1+H)\tau_{\max}}{1-C\tau_{\max}}B_l .
\label{eq:reward-step-explicit}
\end{align}
For the fee components used below, $B_l=1$.
Hence the feedback is $O(\tau_{\max})$ in both coordinates.
The reward coordinate gives
\begin{equation}
\frac{df_k}{dx_l}=\rho^R\,\mathbbm{1}\{k=l\}\sum_j s^*_{jl}\,\frac{\partial\tau_{jl}}{\partial x_l}+O(\tau_{\max}),
\label{eq:rewards-step}
\end{equation}
which confirms that the share-feedback term in~\eqref{eq:reward-decomp} is small.
The price coordinate gives
\begin{equation}
\frac{d\log p_j}{dx_l}=\rho^P\,\mu_{jl}\,\frac{\partial\tau_{jl}}{\partial x_l}+O(\tau_{\max}),
\label{eq:price-step}
\end{equation}
This threshold $\bar\tau=1/C$ is the explicit threshold the theorem promises.

% \begin{remark}[Income effects are first-order, and controlled by the contraction]\label{rem:income-effects}
% Income effects are not negligible.
% Rewards track the fee one-for-one, so $df/dx_l=O(1)$, and the income feedback through $G$ contributes $O(\tau_{\max})$ to the welfare derivative---the \emph{same} order as the price feedback through $E$, not smaller.
% What removes them from the first-order statistic is not that they are higher-order but that the joint contraction~\eqref{eq:contraction-bound} bounds them; this is exactly why the income bound $L$ appears in the threshold $\bar\tau$.
% Treating the reward change itself as $O(\tau_{\max})$---and hence the income feedback as $O(\tau_{\max}^2)$---would be an error: it is the \emph{derivative} $df/dx_l$ that is $O(1)$, not negligible, even though the reward \emph{level} $f_k$ is $O(\tau_{\max})$.
% \end{remark}

% \begin{remark}[Why the feedback vanishes from the statistic]\label{rem:feedback}
% Both feedbacks---price ($E$) and income ($G$)---are $O(\tau_{\max})$ in the welfare \emph{derivative}, and hence $O(\tau_{\max}^2)$ in integrated welfare (an $O(\tau_{\max})$ error in the derivative times a fee of size $O(\tau_{\max})$).
% This is why the equilibrium reallocation of baskets and the income feedback both drop out of the first-order sufficient statistic, which keeps only the mechanical direct effects.
% \end{remark}

\paragraph{Step D: Assembly.}
Substitute the price response~\eqref{eq:price-step} and the reward response~\eqref{eq:rewards-step} into the spending-share envelope formula~\eqref{eq:cohort-roy}:
\[
\frac{1}{\lambda_k}\frac{dV_k}{dx_l}=-\rho^P\sum_j s^*_{jk}\,\mu_{jl}\,\frac{\partial\tau_{jl}}{\partial x_l}+\mathbbm{1}\{k=l\}\,\rho^R\sum_j s^*_{jk}\,\frac{\partial\tau_{jk}}{\partial x_k}+O(\tau_{\max}).
\]
Now translate the cohort spending-weighted sums into observable revenue-weighted moments with the spending identity~\eqref{eq:spending-identity}, $\sum_j s^*_{jk}(\cdot)=\mu_k^{-1}\mathbb{E}_R[\mu_{jk}(\cdot)]+O(\tau_{\max})$:
\begin{equation}
\frac{1}{\lambda_k}\frac{dV_k}{dx_l}=\mu_k^{-1}\!\left(-\rho^P\,\mathbb{E}_R\!\left[\mu_{jk}\mu_{jl}\,\frac{\partial\tau_{jl}}{\partial x_l}\right]+\mathbbm{1}\{k=l\}\,\rho^R\,\mathbb{E}_R\!\left[\mu_{jk}\,\frac{\partial\tau_{jk}}{\partial x_k}\right]\right)+O(\tau_{\max}),
\end{equation}
which is equation~\eqref{eq:sufficient-stat-formula} and proves Theorem~\ref{thm:sufficient-stat}.\qed

\subsection{Proof of Fee Component Formulas}\label{subsec:fee-components}

We derive the fee component decomposition by applying Theorem~\ref{thm:sufficient-stat} to specific counterfactual changes in the fee structure.
We set $\rho^P=\rho^R=1$ throughout for simplicity in writing the expressions.
The welfare effect of any fee change equals the size of the change multiplied by the derivative from Theorem~\ref{thm:sufficient-stat}, so it is enough to track the derivative for each component.

\paragraph{Baseline Fees.}
The baseline fee $\overline{\tau}_l$ shifts every interchange rate in~\eqref{eq:interchange-fee-model} by the same amount, so $\partial\tau_{jl}/\partial\overline{\tau}_l=1$ for all merchants.
Applying Theorem~\ref{thm:sufficient-stat} and using $\mathbb{E}_R[\mu_{jk}]=\mu_k+O(\tau_{\max})$,
\[
A_{kl}\equiv\overline{\tau}_l\,\frac{1}{\lambda_k}\frac{dV_k}{d\overline{\tau}_l}=\overline{\tau}_l\,\mu_k^{-1}\!\left(-\mathbb{E}_R[\mu_{jk}\mu_{jl}]+\mathbbm{1}\{k=l\}\,\mu_k\right).
\]

\paragraph{Sector Adjustments.}
Sector adjustments shift fees only for merchants in the relevant sector, so $\partial\tau_{jl}/\partial\tau_{sl}=\mathbbm{1}\{s_j=s\}$.
Write $\omega_s=\sum_{s_j=s}\omega_j$ for sector $s$'s share of total revenue and $\mu_{sl}=\mathbb{E}_R[\mu_{jl}\mid s_j=s]$ for its average share on card $l$.
The normalization $\sum_s\omega_s\mu_{sl}\tau_{sl}=0$ makes the reward term vanish, because
\[
\mathbb{E}_R[\mu_{jl}\,\tau_{s_j,l}]=\sum_s P_R(s_j=s)\,\mathbb{E}_R[\mu_{jl}\mid s_j=s]\,\tau_{sl}=0.
\]
Only the price (second-moment) term remains, and the total effect of sector adjustments is
\[
\sum_s\tau_{sl}\,\frac{1}{\lambda_k}\frac{dV_k}{d\tau_{sl}}=-\mu_k^{-1}\sum_s\tau_{sl}\,P_R(s_j=s)\,\mathbb{E}_R[\mu_{jk}\mu_{jl}\mid s_j=s].
\]

\paragraph{Negotiated Rates.}
The negotiated discount $\delta_l$ shifts the fee only for large merchants, so $\partial\tau_{jl}/\partial\delta_l=b_j$, which equals one for large merchants and zero otherwise.
Applying Theorem~\ref{thm:sufficient-stat},
\[
B_{kl}\equiv\delta_l\,\frac{1}{\lambda_k}\frac{dV_k}{d\delta_l}=\delta_l\,\mu_k^{-1}\,P_R(b_j=1)\!\left(-\mathbb{E}_R[\mu_{jk}\mu_{jl}\mid b_j=1]+\mathbbm{1}\{k=l\}\,\mathbb{E}_R[\mu_{jk}\mid b_j=1]\right).
\]

\end{document}
